\section{How to Attack}

\begin{rulebox}{General Rule}
The active player is considered the attacker and the inactive player is considered the defender during each combat phase. The attacker's units are called attacking units, and the defender's units are called defending units. For each attack announced at the beginning of the combat phase, the attacker calculates a combat ratio based on each unit's combat strength, which he indexes against a die roll on the Combat Table to resolve the combat.
\end{rulebox}

\ruleslabel{Procedure}

Resolve each combat by executing the following steps in order:

\begin{steps}
  \item The attacker announces which of his Heavy Horse units that are attacking disrupted enemy foot units out of or into clear, slope, or hilltop hexes are conducting a charge. The combat strength of the charging units is doubled.
  \item The attacker halves the combat strengths of all of his units attacking out of a marsh hex, across a ditch hexside or stream hexside, or up a slope hexside, rounding fractions down.
  \item The attacker totals the combat strengths (modified as in steps 1 and 2) of all attacking units participating in the combat into a single combat strength, and adds to it the command rating of any one eligible leader stacked with each attacking unit.
  \item The defender combines the combat strengths of all defending units participating in the combat into a single total combat strength, and adds to it the command rating of any one eligible leader per defending unit that is stacked with that leader.
  \item The defender finds the terrain line on the Combat Table (one of the five lines at the top of the table containing a terrain type and a series of number ratios) that most closely represents the terrain occupied by the defending units. If the defenders occupy more than one terrain type, the terrain type most favorable to the defender is used to resolve the combat. This is the terrain line for the combat.
  \item The attacker expresses the results of steps 3 and 4 as a ratio of attacking units' combat strength to defending units' combat strength, and finds the combat ratio on the terrain line selected in step 5 that most closely approximates this ratio of strength. If the ratio is not exactly equal to one of the combat ratios listed, round the calculated ratio downward to conform to one of the listed ratios.
  \item The attacker rolls one six-sided die and finds the die-roll result in the left-hand column of the Combat Table. Read across the row corresponding to the die roll until reaching the combat-ratio column found in step 6. The entry where the row and column meet is the combat result. Apply all combat results immediately before resolving any other combats. When all combats announced at the beginning of the combat phase have been resolved and all results applied, the combat phase is over.
\end{steps}

\ruleslabel{Cases}

\caseheading{10.1} The effects of terrain on defending units are incorporated into the Combat Table.

There are five terrain lines at the top of the Combat Table. Only the terrain line matching the terrain occupied by the defending units is used to resolve any given combat. The Train/Close---Horse terrain line is used only when defending Heavy or Light Horse units occupy a train or close hex. Similarly, the Train/Woods---Foot terrain line is used only when defending foot units occupy a train or woods hex, and the Close---Foot terrain line is used only when the defending foot units occupy a close hex. If the defender occupies more than one type of terrain in a combined combat, the line for the occupied terrain most favorable to the defender is used. Terrain lines are organized on the Combat Table so that the lower the terrain line, the more favorable it is to the defender.

\caseheading{10.2} Ordered Heavy Horse units can charge disrupted enemy foot units (only).

When one or more ordered Heavy Horse units are attacking one or more disrupted enemy foot units, the attacker can, at his option, announce a charge. Charges can be conducted only if all the attacking Heavy Horse units involved in the combat charge and only if all of the defending units are disrupted foot units. Heavy Horse can charge if friendly foot units are involved in the attack, so long as the other requirements for conducting a charge are met.

Heavy Horse can charge into or out of clear, hilltop, and slope hexes, and can charge across ditch and stream hexsides. However, they can only charge across slope hexsides if they are charging ``upslope'' (for example, from the slope hex forming that slope hexside into a hex that does not contain part of the splash contour of the slope forming that slope hexside). In minimal visibility, the Heavy Horse units must be next to or stacked with an eligible leader at the beginning of the combat phase to announce a charge. See case 13.9 for special charge restrictions placed on Heavy Horse units of demoralized armies. Charging Heavy Horse units have their combat strengths doubled.

\caseheading{10.3} Charging Heavy Horse units that are ordered after the results of their combat have been applied are immediately disrupted. Charging horse units that were disrupted as a result of their combat suffer no further adverse result.

\caseheading{10.4} The command ratings of eligible leaders are added to the combat strength of units with which they are stacked.

A leader is eligible if it is the same color as the unit with which it is stacked. However, the amount added to the unit's combat strength cannot exceed the unit's printed combat strength. If the printed combat strength on the face-up side of the unit is less than the command rating of the eligible leader, only an amount equal to the unit's printed combat strength is added. The command rating of only one eligible leader can be added to a given unit's combat strength, though any number of leaders can add their command ratings to a given combat if they are stacked with different units.
